\section{Background}
\label{sec:background}

Argo floats park and drift at 1000\,m depth between profiling cycles; their
locations are set by currents at that depth, not by experimental design. If those
currents are correlated with upper-ocean temperature, then Argo-based estimates
of space--time-averaged ocean heat content (OHC) can be biased. This is widely
suspected to be negligible, but the magnitude has not been bounded statistically.

Sampling of this kind---locations generated by a process correlated with the
field under study---is \emph{preferential} in the sense of \citet{diggle2010}:
the location process and the field are statistically dependent, so averages formed
from the observations need not match averages over the domain. Its effect on areal
averages and spatial prediction is documented across environmental monitoring and
ecology \citep{gelfand2012,zidek2014,conn2017,dumelle2022}, and the standard
response is a \emph{model-based correction}---a joint model for the locations and
the latent field, a shared Gaussian process in the canonical treatment of
\citet{diggle2010}, whose fit removes the bias. Such corrections require a model
for the sampling mechanism and a covariance estimated from the very data whose
locations are preferential, with the model risk both entail.

We take the complementary, design-based route: rather than correct the bias, we
ask first whether it needs correcting---how large the OHC preferential-sampling
bias can be, in the worst case, if one does nothing. The do-nothing estimator is
the naive area average, which inherits the entire discrepancy between where the
floats sample and the domain being averaged. We bound its bias within a smoothness
class for the OHC field, using only the observed float positions, with the unknown
field entering the guarantee through a single, a-priori-capped amplitude. The
device is the kernel mean embedding of distributions and the maximum mean
discrepancy \citep{smola2007,gretton2012,muandet2017}: the bias is an inner
product against the embedding difference $\mP-\mQ$, and the discrepancy factor is
the MMD between the sampling and reference laws, estimated directly from the
profile positions.

This worst case bounds the estimators actually used. OHC products interpolate the
profiles---by kriging, objective analysis, or optimal interpolation---and
integrate the interpolant \citep{baugh2022}; each is a linear functional of the
observations, of which the naive average is a suboptimal instance and kriging the
optimal one, so the worst-case bound provably upper-bounds the
preferential-sampling bias of any such interpolation estimator
(Proposition~\ref{prop:kb-estimator}). It does not bound nonlinear, model-based
de-biasing, which uses the locations themselves as information; instead it makes
such a correction's necessity assessable. If the worst-case bias already sits
within the interpolation error the products carry, no correction is warranted for
OHC; if it does not, the bound quantifies what is at stake.

\section{A kernel-discrepancy bound on preferential-sampling bias}
\label{sec:kernel}

\subsection{Setup and estimand}
\label{sec:kb-setup}

Fix a measurable space $(\X,\F)$: a bounded space--time domain, a generic point
$x=(s,t)$ carrying spatial location $s$ and time $t$. Two probability measures
live on it. The \emph{reference} (target) measure $\Qb$ is normalized area
$\times$ time---the average an OHC estimate aims at. The \emph{sampling} measure
$\Pb$ is the intensity of the profile locations; it need not equal $\Qb$, and the
gap between them is preferential sampling. We write $X\sim\Qb$ and $Z\sim\Pb$ for
reference and sampled random points.

The OHC field is a fixed function $Y:\X\to\R$ (depth-integrated heat content per
unit area). The estimand is its reference average, a linear functional of the
field \citep{baugh2022},
\[
  \theta \;:=\; \int_\X Y\dd\Qb \;=\; \E_{X\sim\Qb}[Y(X)] ,
\]
and its analogue under the sampling law is the \emph{design mean}
$\theta_{\Pb}:=\int_\X Y\dd\Pb=\E_{Z\sim\Pb}[Y(Z)]$, the expectation of the naive
profile average $N^{-1}\sum_i Y(Z_i)$. The \emph{preferential-sampling bias} is
their difference,
\[
  B \;:=\; \theta_{\Pb}-\theta \;=\; \int_\X Y\dd\Pb-\int_\X Y\dd\Qb ,
\]
the object bounded below. The field enters through a reproducing-kernel smoothness
class.

\begin{assumption}[Field smoothness]
\label{ass:kb}
Let $k:\X\times\X\to\R$ be a symmetric, positive-definite, $\F$-measurable kernel,
bounded by $\sup_{x}k(x,x)=:\kappa_k<\infty$, with reproducing-kernel Hilbert space
$\Hcal$, inner product $\innerHc{\cdot}{\cdot}$ and norm $\normHc{\cdot}$
\citep{aronszajn1950,berlinet2004}. The field lies in $\Hcal$ with bounded norm,
\[
  Y\in\Hcal,\qquad \normHc{Y}\le R,
\]
for a radius $R>0$.
\end{assumption}

Boundedness gives the continuous embedding $\Hcal\hookrightarrow L^2(\rho)$ for any
probability measure $\rho$, with $\|f\|_{L^2(\rho)}\le\sqrt{\kappa_k}\,\normHc{f}$
\citep[eq.~3.13]{muandet2017}, so $\theta$ and $\theta_{\Pb}$ are finite and the
estimand is unchanged. Assumption~\ref{ass:kb} restricts the field to be smooth at
the scale encoded by $k$; its conservative, GLORYS-free calibration and the status
of $R$ are the subject of \S\ref{sec:kb-radius}.

\subsection{The bias as a kernel discrepancy}
\label{sec:kb-representer}

The bias is a difference of expectations of the fixed field under two measures, so
the kernel mean embedding turns it into a single inner product.

\begin{definition}[Kernel mean embedding]
\label{def:kb-embedding}
For a probability measure $\rho$ on $(\X,\F)$ the \emph{kernel mean embedding} is
\[
  \mu_\rho \;:=\; \int_\X k(x,\cdot)\dd\rho(x)\ \in\ \Hcal
\]
\citep[Definition~3.1]{muandet2017}. Write $\mP$ and $\mQ$ for the embeddings of
the sampling and reference measures.
\end{definition}

\begin{lemma}[Existence and the mean property]
\label{lem:kb-mean}
Under Assumption~\ref{ass:kb}, for $\rho\in\{\Pb,\Qb\}$ the embedding $\mu_\rho$ is
a well-defined element of $\Hcal$, and every $f\in\Hcal$ satisfies
\[
  \E_{\rho}[f] \;=\; \innerHc{f}{\mu_\rho}
\]
\citep[Lemma~3.1]{muandet2017,smola2007}.
\end{lemma}

\begin{proposition}[Discrepancy representation of the bias]
\label{prop:kb-representer}
Under Assumption~\ref{ass:kb},
\[
  B \;=\; \innerHc{Y}{\mP-\mQ},
  \qquad
  \normHc{\mP-\mQ} \;=\; \MMDpq ,
\]
where $\MMDpq:=\normHc{\mP-\mQ}$ is the maximum mean discrepancy between the
sampling and reference laws \citep[eq.~3.29]{muandet2017,gretton2012}.
\end{proposition}
\deferproof{kb-representer}{%
\begin{proof}[Proof of Proposition~\ref{prop:kb-representer}]
By Lemma~\ref{lem:kb-mean} applied to $f=Y$,
$\theta_{\Pb}=\innerHc{Y}{\mP}$ and $\theta=\innerHc{Y}{\mQ}$; subtracting and
using linearity of the inner product gives
$B=\theta_{\Pb}-\theta=\innerHc{Y}{\mP-\mQ}$. The norm identity is the definition
of the MMD as the RKHS distance between embeddings \citep[eq.~3.29]{muandet2017}.
\end{proof}%
}

\begin{theorem}[Kernel-discrepancy bound]
\label{thm:kb}
Under Assumption~\ref{ass:kb},
\[
  |B|\;\le\; R\,\MMDpq ,
\]
and the bound is sharp over the smoothness ball:
$\sup_{\normHc{Y}\le R}|B|=R\,\MMDpq$, attained when $\mP\ne\mQ$ at
$Y^{\star}=R\,(\mP-\mQ)/\normHc{\mP-\mQ}$.
\end{theorem}
\deferproof{kb}{%
\begin{proof}[Proof of Theorem~\ref{thm:kb}]
By Proposition~\ref{prop:kb-representer} and Cauchy--Schwarz in $\Hcal$,
$|B|=|\innerHc{Y}{\mP-\mQ}|\le\normHc{Y}\,\normHc{\mP-\mQ}\le R\,\MMDpq$. Equality
in Cauchy--Schwarz holds iff $Y\parallel(\mP-\mQ)$; the constraint $\normHc{Y}\le R$
binds at $Y^{\star}$, for which $\innerHc{Y^{\star}}{\mP-\mQ}=R\,\MMDpq$.
\end{proof}%
}

\begin{remark}[Non-preferential sampling]
\label{rem:kb-zero}
The bound vanishes exactly when $\mP=\mQ$, i.e. $\MMDpq=0$. If $k$ is
characteristic then this occurs iff $\Pb=\Qb$
\citep{sriperumbudur2010}, so the discrepancy is zero precisely for
non-preferential sampling.
\end{remark}

\subsection{A probabilistic bound for a random field}
\label{sec:kb-random}

Theorem~\ref{thm:kb} is deterministic and worst-case: it holds for the single
realized field, uniformly over the smoothness ball, and is attained only when the
field aligns with the embedding difference $\mP-\mQ$. A complementary, average-case
statement models the OHC field as random and asks for the typical magnitude of the
bias, needing only a second moment rather than a norm bound.

\begin{assumption}[Random field]
\label{ass:kb-random}
The field $Y$ is a mean-square random field on $\X$ with mean function
$m(x)=\E[Y(x)]$ and covariance function $c(x,x')=\Cov\big(Y(x),Y(x')\big)$, of
finite second moment: $\int_\X c(x,x)\dd\Pb<\infty$ and $\int_\X c(x,x)\dd\Qb<\infty$.
\end{assumption}

Conditional on the design, the bias $B=\int Y\dd\Pb-\int Y\dd\Qb$ is then a random
variable whose first two moments are kernel discrepancies of the design contrast
$\eta:=\Pb-\Qb$.

\begin{proposition}[Moments of the sampling bias]
\label{prop:kb-moments}
Under Assumption~\ref{ass:kb-random},
\[
  \E[B]=\int_\X m\dd\eta,
  \qquad
  \Var(B)=\iint_{\X\times\X} c(x,x')\dd\eta(x)\dd\eta(x')
         =\MMD[\Hcal_c,\Pb,\Qb]^2 ,
\]
the squared maximum mean discrepancy between $\Pb$ and $\Qb$ under the field's own
covariance kernel $c$, with $\Hcal_c$ its RKHS. If $c=\tau^2 k$ for the smoothness
kernel of Assumption~\ref{ass:kb}, then $\Var(B)=\tau^2\,\MMDpq^2$.
\end{proposition}
\deferproof{kb-moments}{%
\begin{proof}[Proof of Proposition~\ref{prop:kb-moments}]
Conditional on the design $\eta=\Pb-\Qb$ is fixed and $B=\int_\X Y\dd\eta$. The
finite second moment justifies Fubini, so $\E[B]=\int_\X\E[Y]\dd\eta=\int_\X m\dd\eta$
and
\[
  \Var(B)=\E\Big[\big(\!\int_\X(Y-m)\dd\eta\big)^{2}\Big]
  =\iint_{\X\times\X}\Cov\big(Y(x),Y(x')\big)\dd\eta(x)\dd\eta(x')
  =\iint c\dd\eta\dd\eta .
\]
Expanding $\eta=\Pb-\Qb$ gives
$\E_{\Pb\otimes\Pb}[c]-2\E_{\Pb\otimes\Qb}[c]+\E_{\Qb\otimes\Qb}[c]$, which is
$\MMD[\Hcal_c,\Pb,\Qb]^2$ by Proposition~\ref{prop:kb-mmd} applied to the kernel
$c$; for $c=\tau^2 k$ it equals $\tau^2\MMDpq^2$.
\end{proof}%
}

\begin{corollary}[Probabilistic magnitude bound]
\label{cor:kb-cheb}
Under Assumption~\ref{ass:kb-random}, for any $\lambda>0$, with probability at least
$1-\lambda^{-2}$,
\[
  |B|\ \le\ \Big|\int_\X m\dd\eta\Big|+\lambda\,\MMD[\Hcal_c,\Pb,\Qb].
\]
If the mean field is constant, so $\int_\X m\dd\eta=0$, and $c=\tau^2 k$, then $B$ has
mean zero and
\[
  \Pr\big(|B|\ge\lambda\,\tau\,\MMDpq\big)\ \le\ \lambda^{-2},
\]
sharpening to a tail of $2\Phi(-\lambda)$ when $Y$ is Gaussian.
\end{corollary}
\deferproof{kb-cheb}{%
\begin{proof}[Proof of Corollary~\ref{cor:kb-cheb}]
Chebyshev's inequality gives $\Pr(|B-\E B|\ge\lambda\sqrt{\Var B})\le\lambda^{-2}$;
Proposition~\ref{prop:kb-moments} identifies $\sqrt{\Var B}=\MMD[\Hcal_c,\Pb,\Qb]$,
and $|B|\le|\E B|+|B-\E B|$ gives the first display. A constant mean makes
$\int m\dd\eta=0$ (both $\Pb,\Qb$ are probability measures) and $c=\tau^2 k$ gives
$\sqrt{\Var B}=\tau\MMDpq$, the second. If $Y$ is Gaussian then $B$ is Gaussian and
the tail is exactly $2\Phi(-\lambda)$.
\end{proof}%
}

\begin{remark}[Worst case versus typical case]
\label{rem:kb-random}
The two bounds share the design factor $\MMDpq$ and differ only in the field
amplitude: the deterministic radius $R$, over which the bias is maximized, is
replaced by $\lambda\tau$ with $\tau$ the field's marginal standard deviation---a
directly interpretable amplitude bounded a priori by $M/2$ (Popoviciu's
inequality). Crucially $\tau$ is a variance-scale amplitude, so Popoviciu applies
to it, whereas the RKHS radius $R=\normHc{Y}$ is not a variance and Popoviciu does
not bound it ($R\ge\sigma_Y$ in general); the deterministic radius instead needs
the smoothness ceiling $R_{\mathrm{phys}}\asymp M\,\ell_{\min}^{-s}$ of
\S\ref{sec:kb-radius}. The probabilistic amplitude is thus the more directly
boundable of the two. The probabilistic bound is typically far smaller, since
$\Var(B)=\iint c\dd\eta\dd\eta$ measures the field's dispersion \emph{in the single
direction} $\mP-\mQ$---a quadratic form in the field covariance of the type
appearing in Gaussian-process functional variances \citep{baugh2022}---whereas the
worst case assumes the field concentrates there. The mean contribution
$\int m\dd\eta$ is itself the deterministic sampling bias of the climatological
mean field, bounded by $\normHc{m}\MMDpq$ when $m\in\Hcal$, so the mean enters only
through its norm, never through estimation.
\end{remark}

\subsection{A non-stationary convolution kernel}
\label{sec:kb-kernel}

Both bounds leave the kernel $k$ abstract; the estimator needs a concrete choice,
and in the case-study region that choice must be non-stationary. Ocean heat content
across the North Atlantic Current decorrelates over a few tens of kilometres normal
to the meandering front, while along the jet and through the gyre interior it varies
far more slowly, so any single stationary, isotropic kernel compromises a length
scale and orientation that are in truth local. We therefore take $k$ in the
process-convolution tradition of \citet{higdon1998}, whose model of North Atlantic
Ocean temperatures introduced spatially varying convolution kernels, in the
closed-form non-stationary generalization of \citet{paciorek2006} building on
\citet{higdon1999}.

To each location $x\in\X$ attach a symmetric positive-definite \emph{kernel matrix}
$\Sigma(x)$---the local anisotropy ellipse, encoding the correlation range and
orientation at $x$---varying smoothly over the domain. For an isotropic correlation
$\rho_\nu$ valid on $\R^d$ for every $d$, in particular the Mat\'ern correlation of
fixed smoothness $\nu$, the kernel
\[
\begin{gathered}
  k(x,x')
  = |\Sigma(x)|^{1/4}\,|\Sigma(x')|^{1/4}\,
    \Big|\tfrac{\Sigma(x)+\Sigma(x')}{2}\Big|^{-1/2}
    \rho_\nu\!\big(\sqrt{Q(x,x')}\big),\\[2pt]
  Q(x,x')=(x-x')^{\!\top}\Big(\tfrac{\Sigma(x)+\Sigma(x')}{2}\Big)^{-1}(x-x') ,
\end{gathered}
\]
is symmetric and positive-definite \citep[Theorem~1]{paciorek2006}, with $k(x,x)=1$,
hence $\kappa_k=1$. It measures correlation through the Mahalanobis distance $Q$ in
the \emph{averaged} local metric $(\Sigma(x)+\Sigma(x'))/2$, and reduces to the
stationary anisotropic Mat\'ern when $\Sigma(\cdot)$ is constant. It meets
Assumption~\ref{ass:kb}: a bounded PD kernel whose RKHS is norm-equivalent to a
Sobolev space of the order $\nu+d/2$ fixed by $\nu$, while the local anisotropy
$\Sigma(x)$ carries the non-stationarity. On the space--time domain $x=(s,t)$ we
take $\Sigma$ block-diagonal, $\Sigma(x)=\mathrm{diag}\big(\Sigma_{\mathrm s}(s,t),\,\ell_t^2\big)$,
with $\Sigma_{\mathrm s}$ a $2\times2$ spatial anisotropy and $\ell_t$ a temporal
scale.

The anisotropy field $\Sigma(x)$ is the single modelling choice, made without a
reanalysis. The smoothness order $\nu$ governs membership $Y\in\Hcal$ and is fixed
conservatively by fiat: a rougher $\nu$ only enlarges the class, so containment is
monotone in it. The local ranges and orientations in $\Sigma(x)$ are read from
directional Argo variograms in moving windows and floored to a conservative default
wherever the float sampling is too sparse to resolve them, so the fit is trusted
only where the point pattern supports it, using the known positions and no field
values. The length scales in $\Sigma(x)$ are a separate matter from $\nu$: at fixed
smoothness order they leave the class unchanged and act only on the norm and on the
design factor $\MMDhat$, and the latter is \emph{non-monotone} in scale---both a
very smooth and a very rough kernel drive $\MMDhat$ to zero, with a maximum at an
intermediate scale. A shorter scale is therefore not automatically conservative; we
report $\MMDhat$ across a sweep of scales and take its maximum over the plausible
range as the conservative summary. The front geometry or a reanalysis may anchor
$\Sigma(x)$ in a sensitivity analysis, never in the guarantee.

\subsection{Estimating the discrepancy from profile locations}
\label{sec:kb-estimation}

The cap factors into a field radius $R$ (\S\ref{sec:kb-radius}) and a design
factor $\MMDpq$ that depends only on the two measures and the kernel. The latter
is estimable directly from the profile positions, with no intensity model and no
bandwidth. The design factor $k(Z_i,Z_{i'})$ of the non-stationary kernel
(\S\ref{sec:kb-kernel}) is evaluated at the profiles' own locations, so the local
anisotropy enters the discrepancy exactly where the floats sample; the reference
terms below lose the separable closed form and are computed by Monte Carlo over
$\Qb$, carrying an explicit Monte-Carlo error alongside the sampling error.

\begin{proposition}[Closed form of the squared discrepancy]
\label{prop:kb-mmd}
Under Assumption~\ref{ass:kb},
\begin{equation}
  \MMDpq^{2}
  =\E_{\Pb\otimes\Pb}\!\big[k(Z,Z')\big]
    -2\,\E_{\Pb\otimes\Qb}\!\big[k(Z,X)\big]
    +\E_{\Qb\otimes\Qb}\!\big[k(X,X')\big],
  \label{eq:kb-mmd}
\end{equation}
with $Z,Z'\sim\Pb$ and $X,X'\sim\Qb$ independent
\citep[eq.~3.30]{muandet2017}.
\end{proposition}
\deferproof{kb-mmd}{%
\begin{proof}[Proof of Proposition~\ref{prop:kb-mmd}]
Expand $\normHc{\mP-\mQ}^2=\innerHc{\mP}{\mP}-2\innerHc{\mP}{\mQ}+\innerHc{\mQ}{\mQ}$
and use $\innerHc{\mu_\rho}{\mu_{\rho'}}=\iint k(x,x')\dd\rho(x)\dd\rho'(x')$, two
applications of Lemma~\ref{lem:kb-mean}.
\end{proof}%
}

The reference terms are known functions of $\Qb$ and $k$ alone. Writing
\[
  \kbarQ(z):=\int_\X k(z,x)\dd\Qb(x)=\innerHc{k(z,\cdot)}{\mQ},
  \qquad
  \kbbarQ:=\iint_{\X\times\X} k\dd\Qb\dd\Qb=\normHc{\mQ}^2 ,
\]
both are computed by Monte Carlo over $\Qb$ for the non-stationary kernel of
\S\ref{sec:kb-kernel}, without any field product; were $k$ stationary or separable
in the equal-measure coordinates $(\lambda,u,t)=(\text{lon},\sin\phi,\text{time})$
in which $\Qb$ is uniform, $\kbarQ$ would reduce to a product of one-dimensional
integrals and $\kbbarQ$ to a constant. Only the first term of \eqref{eq:kb-mmd}
involves the unknown $\Pb$, and it is estimated from the profiles.

A profile is one Argo cycle; a float contributes many, and its successive cycles
are near one another, so profiles from a single float are not independent. We index
profiles by float and average the kernel over pairs from \emph{distinct} floats, so
that within-float dependence never enters.

\begin{definition}[Between-float discrepancy estimator]
\label{def:kb-hat}
Index the $N$ profiles by float, float $\ell$ contributing $d_\ell$ of them. Let
$\mathcal B$ be the ordered pairs of profiles from distinct floats,
$P_{\mathrm{bf}}=|\mathcal B|=N(N-1)-\sum_\ell d_\ell(d_\ell-1)$. The
\emph{between-float} estimator of \eqref{eq:kb-mmd} is
\begin{equation}
  \MMDhat^{2}
  \;=\;\frac{1}{P_{\mathrm{bf}}}\sum_{(i,i')\in\mathcal B}k(Z_i,Z_{i'})
      \;-\;\frac{2}{N}\sum_{i=1}^{N}\kbarQ(Z_i)
      \;+\;\kbbarQ .
  \label{eq:kb-hat}
\end{equation}
\end{definition}

This is the unbiased MMD U-statistic of \citet[eq.~3.32]{muandet2017} and
\citet{gretton2012} specialized to our design: the reference ($\Qb$) terms are
computed exactly rather than from a second sample, and the sampling ($\Pb$)
U-statistic is restricted to between-float pairs.

\begin{lemma}[Unbiasedness]
\label{lem:kb-unbiased}
If profiles from distinct floats are independent draws with marginal $\Pb$, then
$\E[\MMDhat^{2}]=\MMDpq^{2}$, with no assumption on the joint law of a single
float's profiles.
\end{lemma}
\deferproof{kb-unbiased}{%
\begin{proof}[Proof of Lemma~\ref{lem:kb-unbiased}]
Each ordered pair in $\mathcal B$ is two profiles from distinct floats, independent
with marginal $\Pb$, so $\E[k(Z_i,Z_{i'})]=\E_{\Pb\otimes\Pb}[k]$; averaging over
the $P_{\mathrm{bf}}$ pairs is unbiased for the first term of \eqref{eq:kb-mmd}.
Each $Z_i$ has marginal $\Pb$, so
$\E[\kbarQ(Z_i)]=\iint k\dd\Pb\dd\Qb=\E_{\Pb\otimes\Qb}[k]$, and the second term is
unbiased for it. The third term is constant. Summing matches \eqref{eq:kb-mmd}.
\end{proof}%
}

\paragraph{Uncertainty.}
Both a finite-sample deviation bound and a central limit theorem propagate the
sampling error of $\MMDhat$ into an upper confidence bound for $\MMDpq$, each with
the \emph{float} as the independent unit ($L$ of them), matching the between-float
construction. Both view $\MMDhat^2$ as a U-statistic through the centered kernel
\[
  h_{\Qb}(z,z') := k(z,z')-\kbarQ(z)-\kbarQ(z')+\kbbarQ ,
\]
for which $\E_{\Pb\otimes\Pb}[h_{\Qb}]=\MMDpq^2$, and \eqref{eq:kb-hat} is its
between-float average
\[
  \MMDhat^2=\frac{1}{P_{\mathrm{bf}}}\sum_{(i,i')\in\mathcal B}h_{\Qb}(Z_i,Z_{i'})
\]
(exactly when the $d_\ell$ are equal, up to an $O(1/N)$ reweighting otherwise).

\begin{proposition}[Finite-sample deviation bound]
\label{prop:kb-mcdiarmid}
Under Assumption~\ref{ass:kb} ($|k|\le\kappa_k$) with independent floats, replacing
the profiles of one float changes $\MMDhat^2$ by at most
$c_\ell=4\kappa_k\big[d_\ell(N-d_\ell)/P_{\mathrm{bf}}+d_\ell/N\big]$, so $c_\ell\le8\kappa_k/L$
when the $d_\ell$ are equal. By the method of bounded differences
\citep{mcdiarmid1989,gretton2012}, for every $t>0$
\[
  \Pr\big(\MMDhat^2-\MMDpq^2>t\big)\ \le\ \exp\!\Big(-\tfrac{2t^2}{\sum_\ell c_\ell^2}\Big),
\]
and hence, with probability at least $1-\alpha$,
\[
  \MMDpq\ \le\ \Big(\MMDhat^2+\big(\tfrac12\textstyle\sum_\ell c_\ell^2\,\log\tfrac1\alpha\big)^{1/2}\Big)_{+}^{1/2}.
\]
The bound is distribution-free and valid at every $L$, with no non-degeneracy
assumption; it is the limit used in the reported guarantee.
\end{proposition}
\deferproof{kb-mcdiarmid}{%
\begin{proof}[Proof of Proposition~\ref{prop:kb-mcdiarmid}]
Write $\MMDhat^2=A-B+\kbbarQ$ with $A=P_{\mathrm{bf}}^{-1}\sum_{(i,i')\in\mathcal B}k(Z_i,Z_{i'})$
and $B=\tfrac2N\sum_i\kbarQ(Z_i)$. Since $|k|\le\kappa_k$ gives
$|\kbarQ|\le\kappa_k$, replacing float $\ell$'s $d_\ell$ profiles moves each of its
$\kbarQ(Z_i)$ terms across a range $2\kappa_k$, changing $B$ by at most
$\tfrac2N d_\ell\,2\kappa_k$. Float $\ell$ occurs in $2d_\ell(N-d_\ell)$ of the
$P_{\mathrm{bf}}$ ordered between-float pairs, each $k$-value moving by at most
$2\kappa_k$, changing $A$ by at most $P_{\mathrm{bf}}^{-1}2d_\ell(N-d_\ell)2\kappa_k$;
$\kbbarQ$ is constant. Summing gives $c_\ell$, and the balanced case
$d_\ell(N-d_\ell)/P_{\mathrm{bf}}=1/L=d_\ell/N$ yields $8\kappa_k/L$. McDiarmid's
inequality gives the tail; inverting the one-sided bound at level $\alpha$ and
taking positive part and square root gives the confidence bound.
\end{proof}%
}

\begin{theorem}[Asymptotic normality]
\label{thm:kb-clt}
Suppose the $L$ floats are i.i.d.\ and $\Pb\ne\Qb$ with $k$ characteristic, so the
first projection $h_1(z):=\E_{Z'\sim\Pb}[h_{\Qb}(z,Z')]=\innerHc{k(z,\cdot)}{\mP-\mQ}$
(up to an additive constant) has $\zeta_1:=\Var_{Z\sim\Pb}[h_1(Z)]>0$. Then
\[
  \sqrt{L}\,\big(\MMDhat^2-\MMDpq^2\big)\ \xrightarrow{\ d\ }\ \mathcal N\!\big(0,\,4\zeta_1\big),
  \qquad L\to\infty,
\]
\citep[Theorem~3.6]{muandet2017,gretton2012,serfling1980}, giving the asymptotic
upper limit $\MMDpq^2\le\MMDhat^2+z_{1-\alpha}\,2(\hat\zeta_1/L)^{1/2}$. Under the
null $\Pb=\Qb$ the statistic is degenerate ($\zeta_1=0$) and the limit is instead an
infinite mixture of $\chi^2$ variables at rate $L^{-1}$, so the normal
approximation serves the magnitude, not a test of preferentiality.
\end{theorem}
\deferproof{kb-clt}{%
\begin{proof}[Proof of Theorem~\ref{thm:kb-clt}]
$\MMDhat^2$ is a U-statistic with symmetric kernel $h_{\Qb}$ over the $L$
independent floats. By Lemma~\ref{lem:kb-mean},
$h_1(z)=\innerHc{k(z,\cdot)}{\mP}-\kbarQ(z)+\text{const}=\innerHc{k(z,\cdot)}{\mP-\mQ}+\text{const}$,
whose $\Pb$-variance $\zeta_1$ is positive iff $\mP\ne\mQ$, i.e.\ (for characteristic
$k$) iff $\Pb\ne\Qb$. When $\zeta_1>0$ the Hoeffding decomposition is
non-degenerate and the U-statistic central limit theorem gives
$\sqrt L(\MMDhat^2-\MMDpq^2)\to\mathcal N(0,4\zeta_1)$; when $\zeta_1=0$ the leading
term vanishes and the second-order term gives the $\chi^2$-mixture limit at rate
$L^{-1}$.
\end{proof}%
}

\subsection{Any profile-based estimator}
\label{sec:kb-estimators}

Scientists rarely use the naive profile mean; they interpolate the profiles---by
kriging or objective analysis---and integrate the interpolant. The bound extends to
any such estimator, because every linear functional of a field in $\Hcal$ has a
representer.

\begin{proposition}[Worst-case error of a linear estimator]
\label{prop:kb-estimator}
Let $\hat\theta_\omega=\sum_{i=1}^{N}\omega_i\,Y(Z_i)$ be a profile-based linear
estimator with $\sum_i\omega_i=1$. Under Assumption~\ref{ass:kb},
\[
  \hat\theta_\omega-\theta=\innerHc{Y}{\eta_\omega},
  \qquad
  \eta_\omega:=\sum_{i=1}^{N}\omega_i\,k(Z_i,\cdot)-\mQ,
\]
and $\sup_{\normHc{Y}\le R}|\hat\theta_\omega-\theta|=R\,\normHc{\eta_\omega}$. The
weights minimizing $\normHc{\eta_\omega}$ are the kernel (Bayes--Hermite) quadrature
weights---equivalently the kriging predictor of $\theta$---and the minimum is the
associated power function \citep{ohagan1991,micchelli1977,donoho1994}.
\end{proposition}
\deferproof{kb-estimator}{%
\begin{proof}[Proof of Proposition~\ref{prop:kb-estimator}]
By the reproducing property $Y(Z_i)=\innerHc{Y}{k(Z_i,\cdot)}$ and
$\theta=\innerHc{Y}{\mQ}$ (Lemma~\ref{lem:kb-mean}), so
$\hat\theta_\omega-\theta=\innerHc{Y}{\sum_i\omega_i k(Z_i,\cdot)-\mQ}$; the
constraint $\sum_i\omega_i=1$ ensures exactness on constants. Cauchy--Schwarz over
the ball gives $R\normHc{\eta_\omega}$ as in Theorem~\ref{thm:kb}. Minimizing
$\normHc{\eta_\omega}^2$ subject to $\sum_i\omega_i=1$ is the kernel-quadrature /
kriging problem.
\end{proof}%
}

Equal weights $\omega_i=1/N$ make $\eta_\omega$ the empirical form of $\mP-\mQ$,
whose worst-case error is the design bias bounded by $R\,\MMDpq$. Kriging minimizes
$\normHc{\eta_\omega}$, so $\normHc{\eta_{\mathrm{kr}}}\le\normHc{\eta_{1/N}}$: the
naive average is the worst member of the linear family and the optimal (kriging)
member is the best, hence the design-bias bound provably upper-bounds the
preferential-sampling bias of any interpolation estimator. Their radius-free ratio,
\[
  \frac{\sup_{\normHc{Y}\le R}|B|}{\sup_{\normHc{Y}\le R}|\hat\theta_{\mathrm{kr}}-\theta|}
  =\frac{\normHc{\eta_{1/N}}}{\normHc{\eta_{\mathrm{kr}}}}\ \ge\ 1 ,
\]
computable from positions and kernel alone, measures how much ignoring preferential
sampling costs beyond the interpolation error the optimal estimator incurs anyway.

\subsection{The field radius and how the bound is reported}
\label{sec:kb-radius}

Of the two factors in Theorem~\ref{thm:kb}, only $R=\normHc{Y}$ is not fixed by the
profile locations, and it cannot be identified from finitely many observations: the
minimum-norm interpolant $\hat Y$ through the readings has
$\normHc{\hat Y}=(\mathbf y^\top K^{-1}\mathbf y)^{1/2}\le\normHc{Y}$
($K=[k(Z_i,Z_{i'})]$, $\mathbf y=(Y(Z_i))_i$), a \emph{lower} bound---the unsafe
direction for a cap---while point evaluations place no upper bound on the norm. A
finite bound on a linear functional over a function class exists only when the class
is bounded, the minimax error being otherwise infinite
\citep{donoho1994,micchelli1977}. The kernel bound does not remove this irreducible
assumption but isolates it in a single interpretable scalar $R$, an
amplitude-of-roughness in OHC units, capped a priori and never estimated.

\paragraph{A guaranteed, physically ceilinged bound.}
A finite $R$ follows from two declared physical facts, neither using a reanalysis
field. OHC is bounded---seawater temperature lies in a known range, so
$Y\in[a,b]$ over the domain with $M:=b-a$---and carries no structure below a minimum
feature scale $\ell_{\min}$. For a kernel whose RKHS embeds in a Sobolev space of
order $s$, these bound the norm, $\normHc{Y}\le R_{\mathrm{phys}}\asymp M\,\ell_{\min}^{-s}<\infty$,
so $|B|\le R_{\mathrm{phys}}\MMDhat$ is a true guarantee consulting neither an
external field nor the OHC values---loose but rigorous. The two inputs are tied to
the kernel: $\ell_{\min}$ is at once the shortest scale in the smoothness class and
the quantity that keeps $R$ finite, since $\normHc{Y}\to\infty$ as the class is
refined to zero scale. The conservative kernel therefore targets the minimum
\emph{physical} scale, not zero.

\paragraph{What the bias must beat: prediction error.}
The scientifically meaningful question is not the raw size of the bias but whether
it is large relative to the interpolation error already tolerated in OHC products.
We report the breakdown budget $R^{\star}(\mathcal T)=\mathcal T/\MMDhat$, the
roughest the field could be before the bias reaches a benchmark $\mathcal T$, taken
to be the mapping or kriging error the product itself carries; the bias exceeds it
only if $\normHc{Y}>R^{\star}(\mathcal T)$, a parameter-free, falsifiable statement.
Because that prediction error is itself $R\,\normHc{\eta_{\mathrm{kr}}}$
(Proposition~\ref{prop:kb-estimator}), the radius cancels in the ratio of
\S\ref{sec:kb-estimators}, and the comparison of worst-case bias to worst-case
interpolation error needs no radius and no reanalysis. For a plausible (rather than
worst-case) magnitude, $R$ may be anchored in sensitivity analysis by the
interpolant floor $\normHc{\hat Y}$ and by a reanalysis over the full domain; both
are lower anchors to which a margin is added, and neither enters the guarantee.

\begin{remark}[What enters the guarantee]
\label{rem:kb-ledger}
The bound $R\,\MMDhat$ uses only the profile locations, through $\MMDhat$ of
\eqref{eq:kb-hat}, and the field radius $R$; the radius is capped a priori by
$R_{\mathrm{phys}}$ or left free in the breakdown budget, never estimated from the
field. A reanalysis enters only the kernel scale and a sensitivity anchor for $R$,
never a boxed guarantee. The standing assumptions are $Y\in\Hcal$ with
$\normHc{Y}\le R$ (Assumption~\ref{ass:kb}) and between-float independence for the
unbiasedness of $\MMDhat$ (Lemma~\ref{lem:kb-unbiased}). The discrepancy is a single
functional estimated at all scales at once, so no partition of the domain, and no
resolution choice, is required.
\end{remark}

\subsection{A worked example}
\label{sec:kb-example}

We evaluate the estimator on a small synthetic design in the case-study box
($\lambda\in[-68,-62]$, $\phi\in[36,40]$, $t\in[2020,2023)$) to exhibit the
arithmetic; the reference terms are exact and only the profile positions enter the
design factor $\MMDhat$.

\paragraph{Coordinates, kernel, reference terms.}
Map each equal-measure coordinate to the unit interval,
$a=\tfrac{\lambda+68}{6}$, $b=\tfrac{\sin\phi-\sin36^\circ}{\sin40^\circ-\sin36^\circ}$,
$c=\tfrac{t-2020}{3}$, so $\Qb$ is uniform on $[0,1]^3$. Take the separable
exponential (Mat\'ern-$\tfrac12$) correlation kernel $k(x,x')=\prod_{d}e^{-|x_d-x'_d|/\ell_d}$
over $d\in\{a,b,c\}$, with $\ell_a=\ell_b=0.3$, $\ell_c=0.2$; here $k(x,x)=1$, so
the kernel carries correlation shape only and the amplitude resides in $R$. Since
$k$ is a product and $\Qb$ product-uniform, the reference integrals factor through
$g(z;\ell)=\int_0^1 e^{-|z-x|/\ell}\dd x=\ell(2-e^{-z/\ell}-e^{-(1-z)/\ell})$ and
$G(\ell)=\iint_{[0,1]^2}e^{-|x-x'|/\ell}=2\ell-2\ell^2(1-e^{-1/\ell})$, via
$\kbarQ(a,b,c)=g(a;\ell_a)g(b;\ell_b)g(c;\ell_c)$ and $\kbbarQ=G(\ell_a)G(\ell_b)G(\ell_c)$.
Numerically $G(0.3)=0.426$, $G(0.2)=0.321$, so $\kbbarQ=0.058$.

\paragraph{Design and computation.}
Seven profiles from three floats, floats $A,B$ clustered in one corner and float
$C$ apart (Figure~\ref{fig:kb-toy}), with their exact $\kbarQ(Z_i)$:

\begin{center}
\begin{tabular}{llrrrr}
\hline
float & profile & $a$ & $b$ & $c$ & $\kbarQ(Z_i)$\\
\hline
$A$ & $A_1$ & 0.20 & 0.30 & 0.10 & 0.0541\\
$A$ & $A_2$ & 0.24 & 0.33 & 0.20 & 0.0667\\
$A$ & $A_3$ & 0.28 & 0.28 & 0.30 & 0.0723\\
$B$ & $B_1$ & 0.22 & 0.40 & 0.15 & 0.0630\\
$B$ & $B_2$ & 0.26 & 0.38 & 0.55 & 0.0784\\
$C$ & $C_1$ & 0.75 & 0.70 & 0.50 & 0.0752\\
$C$ & $C_2$ & 0.80 & 0.65 & 0.80 & 0.0648\\
\hline
\end{tabular}
\end{center}

\begin{figure}[h]
\centering
\begin{tikzpicture}[
  markA/.style={circle,fill=black,minimum size=4pt,inner sep=0},
  markB/.style={rectangle,fill=black,minimum size=4pt,inner sep=0},
  axis/.style={->,gray}]
\begin{scope}
  \draw[gray!50] (0,0) rectangle (4,4);
  \draw[axis] (0,0) -- (4.5,0) node[right,black]{\small $a$ (lon)};
  \draw[axis] (0,0) -- (0,4.5) node[above,black]{\small $b$ ($\sin\phi$)};
  \foreach \t in {0,0.5,1}{%
    \draw (\t*4,-0.08)--(\t*4,0.08) node[below=1pt,black]{\scriptsize \t};
    \draw (-0.08,\t*4)--(0.08,\t*4) node[left=1pt,black]{\scriptsize \t};}
  \draw[dashed,gray] (0.6,0.9) rectangle (1.3,1.8);
  \node[markA] at (0.8,1.2){}; \node[markA] at (0.96,1.32){}; \node[markA] at (1.12,1.12){};
  \node[markB] at (0.88,1.6){}; \node[markB] at (1.04,1.52){};
  \draw[thick] (2.91,2.71)--(3.09,2.89) (2.91,2.89)--(3.09,2.71);
  \draw[thick] (3.11,2.51)--(3.29,2.69) (3.11,2.69)--(3.29,2.51);
  \node at (2,-0.95){\small space};
\end{scope}
\begin{scope}[xshift=7cm]
  \draw[axis] (0,0) -- (4.5,0) node[right,black]{\small $c$ (time)};
  \foreach \t in {0,0.5,1}{\draw (\t*4,-0.08)--(\t*4,0.08) node[below=1pt,black]{\scriptsize \t};}
  \foreach \r/\lab in {3/A,2/B,1/C}{%
    \node[left,black] at (-0.1,\r){\small $\lab$};
    \draw[gray!30] (0,\r)--(4,\r);}
  \node[markA] at (0.4,3){}; \node[markA] at (0.8,3){}; \node[markA] at (1.2,3){};
  \node[markB] at (0.6,2){}; \node[markB] at (2.2,2){};
  \draw[thick] (1.91,0.91)--(2.09,1.09) (1.91,1.09)--(2.09,0.91);
  \draw[thick] (3.11,0.91)--(3.29,1.09) (3.11,1.09)--(3.29,0.91);
  \node at (2,-0.95){\small time};
\end{scope}
\end{tikzpicture}
\caption{The synthetic design of \S\ref{sec:kb-example}. \emph{Left:} profile
locations in the normalized spatial coordinates $(a,b)$; floats $A$ (filled
circles) and $B$ (squares) cluster in one corner (dashed box) while float $C$
(crosses) samples elsewhere. \emph{Right:} the sampling time $c$ of each profile,
by float. The design factor $\MMDhat$ depends only on this point pattern.}
\label{fig:kb-toy}
\end{figure}

Here $N=7$, float sizes $(3,2,2)$, and $P_{\mathrm{bf}}=42-10=32$ ordered pairs
($16$ unordered). The reference contributions are
$\tfrac{2}{N}\sum_i\kbarQ(Z_i)=\tfrac27(0.4745)=0.136$ and $\kbbarQ=0.058$; the
first term of \eqref{eq:kb-mmd} is the average of $k$ over the $16$ between-float
pairs (within-float pairs excluded), which is $1.884/16=0.118$. Combining as in
\eqref{eq:kb-hat},
\[
  \MMDhat^{2}
  =\underbrace{0.118}_{\hat\E_{\Pb\otimes\Pb}}
  -\underbrace{0.136}_{2\hat\E_{\Pb\otimes\Qb}}
  +\underbrace{0.058}_{\E_{\Qb\otimes\Qb}}
  =0.040,
  \qquad
  \MMDhat=0.201 .
\]
The ordering $\hat\E_{\Pb\otimes\Pb}>\hat\E_{\Pb\otimes\Qb}>\E_{\Qb\otimes\Qb}$ is
the signature of preferential sampling: pairs from the clustered design are closer,
hence higher-$k$, than design-versus-reference or reference pairs; a design spread
like $\Qb$ would collapse the three and drive $\MMDhat\to0$.

\paragraph{The bound.}
By Theorem~\ref{thm:kb}, $|B|\le R\,\MMDhat=0.201\,R$, the positions fixing the
slope and the field budget $R$ the scale. Reported as in \S\ref{sec:kb-radius}: a
breakdown budget $R^{\star}(\mathcal T)=\mathcal T/\MMDhat$, so a tolerance
$\mathcal T=0.05\,\mathrm{GJ/m^2}$ is exceeded only if the field's $\Hcal$-norm
exceeds $0.25$; or a curve $0.201\,R$ up to the physical ceiling
$R_{\mathrm{phys}}$. None of it consults an external field.

\appendix
\section{Proofs}
\label{app:proofs}
\printdeferredproofs
